\chapter{Biography}\label{app:biography}

\begin{BlindRegion}

\newcommand{\tikzinterlines}[6][8]{%
  \path let
    \p1=(#2),
    \p2=(#3),
    \p3=(#4),
    \p4=(#5),
    \n{x12}={(\x1-\x2)/#1},
    \n{y12}={(\y1-\y2)/#1},
    \n{x34}={(\x3-\x4)/#1},
    \n{y34}={(\y3-\y4)/#1},
    \n{xy12}={\x1/#1*\y2/#1-\y1/#1*\x2/#1},
    \n{xy34}={\x3/#1*\y4/#1-\y3/#1*\x4/#1},
    \n{x}={(\n{xy12}*\n{x34}-\n{x12}*\n{xy34})/(\n{x12}*\n{y34}-\n{y12}*\n{x34})},
    \n{y}={(\n{xy12}*\n{y34}-\n{y12}*\n{xy34})/(\n{x12}*\n{y34}-\n{y12}*\n{x34})}
  in
    coordinate (#6) at (#1*\n{x},#1*\n{y});
}

\newcommand\tikzhyperbola[7][thick]{%
  \draw [#1,rotate={#2},shift={#3}]
    plot [variable=\t,samples=100,smooth,domain=#6:#7] ({#4*sec(\t)},{#5*tan(\t)});
  \draw [#1,rotate={#2},shift={#3}]
    plot [variable=\t,samples=100,smooth,domain=#6:#7] ({-#4*sec(\t)},{#5*tan(\t)});
}

\newcommand{\tikztantocircle}[5]{%
  \path let
    \p0=(#1),
    \p{r}=(#2,0),
    \n{r}={\x{r}},
    \p{a}=(#3),
    \n{dx}={\x{a}-\x0},
    \n{dy}={\y{a}-\y0},
    \n{theta0}={atan2(\n{dy},\n{dx})},
    \n{theta}={acos(\n{r}/veclen(\n{dx},\n{dy})},
    \n{x3}={\x0+\n{r}*cos(\n{theta0}+\n{theta})},
    \n{y3}={\y0+\n{r}*sin(\n{theta0}+\n{theta})},
    \n{x4}={\x0+\n{r}*cos(\n{theta0}-\n{theta})},
    \n{y4}={\y0+\n{r}*sin(\n{theta0}-\n{theta})}
  in
    coordinate (#4) at (\n{x3},\n{y3})
    coordinate (#5) at (\n{x4},\n{y4});
}

\newsavebox{\BrianchonPicture}
\sbox{\BrianchonPicture}{
  \begin{tikzpicture}[x=3.8pt,y=3.8pt]
    \def\a{12}
    \def\b{9}
    \def\Rot{70}
    \def\Sh{(0,0)}
    \tikzhyperbola{\Rot}{\Sh}{\a}{\b}{-50}{63}
    \foreach \i/\y/\ty in {1/15/-15,2/3/-15,3/-6/15,4/-8/15,5/2/-15,6/16/-15} {
      \pgfmathsetmacro{\s}{\i <= 3 ? -1 : 1}
      \pgfmathsetmacro{\x}{\s*sqrt(\a*\a*(1+\y*\y/(\b*\b)))}
      \pgfmathsetmacro{\tx}{(1+(\y*\ty)/(\b*\b))/(\x/(\a*\a))}
      \coordinate (r\i) at (\x,\y);
      \coordinate (rt\i) at (\tx,\ty);
    }
    \foreach \i in {1,...,6} {
      \coordinate (\i) at ([rotate={\Rot},shift={\Sh}]r\i);
      \coordinate (t\i) at ([rotate={\Rot},shift={\Sh}]rt\i);
    }
    \foreach \i/\j in {1/2,2/3,3/4,4/5,5/6,6/1} {
      \tikzinterlines{\i}{t\i}{\j}{t\j}{x\i}
    }
    \tikzinterlines{x1}{x4}{x2}{x5}{x7}

    \foreach \i/\j in {1/x6,x1/x2,3/x3,x3/4,x4/x5,6/x6}
      \path[draw=SteelBlue,thick] (\i) -- (\j);

    \foreach \i/\j in {x1/x6,x1/x2,x2/x3,x3/x4,x4/x5,x5/x6}
      \path[draw=SteelBlue,very thick] (\i) -- (\j);

    \foreach \i/\j in {1/4,2/5,3/6}
      \path[draw=Crimson,very thick] (x\i) -- (x\j);

    \foreach \i in {1,...,6} {
      \path[draw=none,fill=SteelBlue] (\i) circle (2pt);
      \path[draw=black,thick,fill=white] (x\i) circle (2pt);
    }

    \path[draw=black,thick,fill=Crimson] (x7) circle (2pt);
  \end{tikzpicture}
}

\newsavebox{\PascalPicture}
\sbox{\PascalPicture}{
  \tikz[x=5pt,y=5pt]{
    \path[draw=black,thick,rotate=70] (0,0) ellipse (12 and 8);

    \foreach \i/\a in {1/10, 2/30, 3/60, 4/100, 5/150, 6/210} {
      \coordinate (r\i) at (\a:12 and 8);
      \coordinate (\i) at ([rotate=70]r\i);
    }

    \tikzinterlines{1}{2}{4}{5}{x14}
    \tikzinterlines{2}{3}{5}{6}{x25}
    \tikzinterlines{3}{4}{6}{1}{x36}
    \foreach \i/\j in {1/4,2/5,3/6}
      \foreach \k in {\i,\j}
        \path[draw=SteelBlue,thick] (\k) -- (x\i\j);
    \path[draw=Crimson,very thick] (x25) -- (x36);
    \path[draw=SteelBlue,very thick] (1) -- (2) -- (3) -- (4) -- (5) -- (6) -- cycle;

    \foreach \i in {1,2,3,4,5,6,x14,x25,x36}
      \path[draw=black,thick,fill=white] (\i) circle (2pt);
  }
}

\newsavebox{\DesarguesPicture}
\sbox{\DesarguesPicture}{
  \begin{tikzpicture}[x=5pt,y=5pt,baseline=(x2)]
    \coordinate (o) at (0,0);
    \foreach \v/\a/\r/\s in {A/85/24/0.4,B/110/30/0.5,C/150/23/0.7}
      \draw[thick,SteelBlue] (o) -- coordinate[pos=\s] (\v1) +(\a:\r) coordinate (\v);

    \tikzinterlines{A}{B}{A1}{B1}{x1}
    \tikzinterlines{B}{C}{B1}{C1}{x2}
    \tikzinterlines{A}{C}{A1}{C1}{x3}

    \draw[SlateGray,thick] (A) -- (x1) (A1) -- (x1);
    \draw[SlateGray,thick] (B) -- (x2) (B1) -- (x2);
    \draw[SlateGray,thick] (A) -- (x3) (A1) -- (x3);

    \draw[black,very thick,fill=Olive,fill opacity=0.75] (A) -- (B) -- (C) -- cycle;
    \draw[black,very thick,fill=OliveDrab,fill opacity=0.75] (A1) -- (B1) -- (C1) -- cycle;

    \draw[Crimson,very thick] (x1) -- (x2);

    \foreach \c/\s in {{white/x1,x2,x3},{SteelBlue/o}}
      \foreach \i in \s
        \path[draw=black,thick,fill=\c] (\i) circle (2pt);
  \end{tikzpicture}
}

\newsavebox{\MongePicture}
\sbox{\MongePicture}{
  \begin{tikzpicture}[x=5pt,y=5pt,baseline=(AB)]
    \pgfmathsetmacro{\rA}{8}
    \pgfmathsetmacro{\rB}{5}
    \pgfmathsetmacro{\rC}{3}
    \pgfmathsetmacro{\rAB}{\rA/(\rA-\rB)}
    \pgfmathsetmacro{\iAC}{\rC/\rA}
    \pgfmathsetmacro{\rBC}{\rB/(\rB-\rC)}

    \draw[very thick,black] (100:32) coordinate (A) circle (\rA);
    \draw[very thick,black] (75:20) coordinate (B) circle (\rB);
    \coordinate (AB) at ($(A)!\rAB!(B)$);
    \path let \p1=(AB),\p2=(A) in
      coordinate (AC) at (\x2-1,\y1)
      coordinate (C) at ($(AC)!\iAC!(A)$);
    \coordinate (BC) at ($(B)!\rBC!(C)$);
    \draw[very thick,black] (C) circle (\rC);

    \foreach \i/\j in {A/B,A/C,B/C} {
      \tikztantocircle{\i}{\csname r\i\endcsname}{\i\j}{\i\j\i1}{\i\j\i2}
      \tikztantocircle{\j}{\csname r\j\endcsname}{\i\j}{\i\j\j1}{\i\j\j2}
    }

    \foreach \c/\i/\s in {{SteelBlue/AB/ABA1,ABA2},{SeaGreen/AC/ACA1,ACA2},{Orchid/BC/BCB1,BCB2}}
      \foreach \j in \s
        \draw[\c,thick] (\i) -- (\j);

    \draw[Crimson,very thick] (AB) -- (BC);

    \foreach \c/\s in {{SteelBlue/ABA1,ABA2,ABB1,ABB2},{SeaGreen/ACA1,ACA2,ACC1,ACC2},{Orchid/BCB1,BCB2,BCC1,BCC2},{Crimson/AB,AC,BC}}
      \foreach \i in \s
        \path[draw=black,thick,fill=\c] (\i) circle (2pt);
  \end{tikzpicture}
}

\begin{wrapfigure}{l}{1.5in}
\resizebox{\linewidth}{!}{\usebox{\BrianchonPicture}} \\
\resizebox{\linewidth}{!}{\usebox{\PascalPicture}}
\end{wrapfigure}
\paragraph{\InsertStudentName}
Brianchon's theorem~\cite{hilbert1932ansgeo} is the dual of Pascal's theorem.
It states that, given a hexagon circumscribed on a conic section, the three lines joining opposite polygon vertices (polygon diagonals) meet in a single point.
\par
Pascal's theorem is the dual of Brianchon's theorem,
discovered by Blaise Pascal in 1640 when he was just 16 years old.
It states that, given a hexagon inscribed in a conic section, the three pairs of the continuations of opposite sides meet on a straight line, called the Pascal line.

\paragraph{\InsertSupervisor}
Desargues' theorem is self-dual.
It states that, if the three straight lines joining the corresponding vertices of two triangles all meet in a point (the perspector),
then the three intersections of pairs of corresponding sides lie on a straight line (the perspectrix).
Equivalently, if two triangles are perspective from a point, they are perspective from a line.
\par
Monge's Theorem~\cite{monge1827desgeo} states that, if we draw three circles in the plane, none of which lies completely inside another, and the common external tangent lines for each pair,
then points of intersection of the three pairs of tangent lines lie on a straight line.
\par
Two books --- ``G\'{e}om\'{e}trie Descriptive'' and ``Anschauliche Geometrie'', the Chinese editions --- influenced me the most on my interest in geometry when I was a teenager.
Over the years, these masterpieces are still so enjoyable to read.

\paragraph{Remarks on Ti\emph{k}Z}
{\tt\textbackslash path} {\tt plot} takes parametric equations, drawing most types of curves.
{\tt\textbackslash pgfmath} and {\tt\textbackslash path} {\tt let} enable us to write very complex mathematics.
However, with 32-bit fixed-point arithmetic, the precision and range are limited.

\begin{figure}[b]
\raggedleft
\resizebox{1.44in}{!}{\usebox{\DesarguesPicture}}
\hspace{2em}
\resizebox{1.8in}{!}{\usebox{\MongePicture}}
\end{figure}

\cleardoublepage\pagestyle{empty}
\def\RuledSurf{
\foreach \i in {0,...,\m} {
  \pgfmathsetmacro{\d}{\i/\m*360+5}
  \coordinate (B\i) at ($(b)+(\d:{\bela} and {\belb})$);
  \coordinate (T\i) at ($(t)+(\d:{\tela} and {\telb})$);
}
\draw[very thick,Crimson] (b) ellipse ({\bela} and {\belb});
\foreach \i in {0,...,\m} {
  \pgfmathsetmacro{\j}{Mod(\i+\f,\m)}
  \pgfmathsetmacro{\k}{Mod(\i-\f,\m)}
  \draw[Plum] (T\i) -- (B\j) coordinate (C\i);
  \draw[Tan] (T\i) -- (B\k)  coordinate (D\i);
}
\draw[very thick,Teal] (t) ellipse ({\tela} and {\telb});

\foreach \i/\j/\k in {5/17/1,14/27/3,23/37/5} {
  \tikzinterlines[16]{T\i}{C\i}{T\j}{D\j}{x\k}
}
\foreach \i/\j/\k in {17/14/2,27/23/4,37/5/6} {
  \tikzinterlines[16]{T\i}{D\i}{T\j}{C\j}{x\k}
}

\foreach \i/\j in {1/2,2/3,3/4,4/5,5/6,6/1}
  \draw[very thick,Coral] (x\i) -- (x\j);

\tikzinterlines[20]{x1}{x4}{x2}{x5}{xx}
\draw[very thick,Crimson] (x1) -- (x4) (x2) -- (x5) (x6) -- (xx);

\foreach \c/\s in {{white/1,2,3,4,5,6},{Crimson/x}}
  \foreach \i in \s
    \path[draw=black,thick,fill=\c] (x\i) circle (1.5pt);
}
\begin{figure}[p]
\centering
\begin{tikzpicture}[x=5pt,y=5pt]
\def\m{70}\def\f{22}
\def\tela{20}\def\telb{6}
\def\bela{10}\def\belb{3}
\coordinate (b) at (0,0);
\coordinate (t) at (-8,100);
\RuledSurf
\def\m{70}\def\f{22}
\def\tela{20}\def\telb{20}
\def\bela{10}\def\belb{10}
\coordinate (b) at (38,40);
\coordinate (t) at (30,40);
\RuledSurf
\end{tikzpicture}
\end{figure}

\end{BlindRegion}

\cleardoublepage
%%
%% END OF THESIS
%%

\endinput
